\chapter{原生智能体与结构反馈闭环}\label{ch:loop}
\section{闭环应该更新什么}
完整问题可以写成一个联合模型：
\begin{equation}
 p(\theta,\alpha,\eta,\beta\mid y,M)
 \propto p(y\mid\theta,\alpha,\eta,\beta)\,
 p(\theta\mid M)p(\alpha)p(\eta)p(\beta).
\end{equation}
$M$ 表示组成、化学连接或候选结构类型。观测 $y$ 固定；闭环更新的是对结构、几何、背景或晶畴的解释。结构因子预测应进入前向模型或作为明确的先验，不能替换原始像素后再把替换结果当作另一份实验。

0.8.0 的可执行实现分成两层。数值层提供有物理定义的求解器、导数、似然和诊断；智能体层规定下一步计算、数据角色、预算和输出。宿主大模型可以提出一个有类型的操作请求，软件只执行注册过的算子。模型生成的文字不会被当作晶体测量，也不会直接作为 Python 或 shell 代码运行。

\section{三种数据角色}
将独立证据分为 $\mathcal W$、$\mathcal S$、$\mathcal A$，分别用于拟合、选择和审计。对第 $m$ 个候选：
\begin{equation}
 \widehat\theta_m=\argmin_\theta L_{\mathcal W}(\theta;M_m),\qquad
 m^*=\argmin_m L_{\mathcal S}(\widehat\theta_m;M_m).
\end{equation}
选定候选后可以用 $\mathcal W\cup\mathcal S$ 再精修一次，随后只评价一次 $\mathcal A$。审计结果写入最终报告，不据此再自动调参。这是本版控制器实际执行的顺序。

不能随意逐行随机分割。等价反射、Friedel 对、同一个像素组的多个分量以及显式声明的共享测量应处于同一个连通分量。代码以 Miller 关系和重叠连接生成图，再按连通分量分配角色；不以观测符号或信噪比选择数据。已有的预留标记会被保留，若标记切开同一证据分量则报错。

这仍然是给定预处理模型下的验证。若仪器几何、轮廓或吸收已经利用了所有数据，强度层的保留集不能宣称整个原始图像处理也完全独立。严格的实验比较还应在外层按扫描块或独立采集划分，并重新执行相应标定。

\begin{figure}[!htbp]
\centering
\begin{tikzpicture}[>=Latex,node distance=8mm and 11mm,
 box/.style={draw=ink!50,fill=pale,rounded corners=2pt,align=center,text width=35mm,minimum height=12mm,font=\small}]
\node[box] (e) {关联证据分组\\等价、Friedel、共像素};
\node[box,below left=of e] (w) {work\\生成与拟合};
\node[box,below=of e] (s) {selection\\选择候选};
\node[box,below right=of e] (a) {audit\\封存};
\node[box,below=18mm of s] (m) {唯一选定模型\\可用 work+selection 再拟合};
\node[box,below=of m] (f) {最终报告\\一次独立评价};
\draw[->] (e)--(w);\draw[->] (e)--(s);\draw[->] (e)--(a);
\draw[->] (w)--node[above,font=\scriptsize]{候选}(s);
\draw[->] (s)--(m);\draw[->] (m)--(f);
\draw[->,accent,thick] (a.east)--++(7mm,0)|-(f.east);
\draw[->,ink!55] (s.south west)--++(-6mm,-8mm)-|(w.south);
\end{tikzpicture}
\caption{控制器的数据职责。候选之间可以在工作与选择证据上比较；审计分支只有通向最终评价的路径。改变审计数据的数值不应改变已经选择的结构。}
\end{figure}

\section{从隔离反射启动，再使用全部区域}
新的 \code{solve-observations} 流程以压缩高斯像素组为输入。初始化时只选几何上隔离、属于参考晶畴的单分量组，不按强度阈值筛选。若该组的压缩矩阵是一列 $r$，则它给出
\begin{equation}
 \widehat I=\frac{r^{\mathsf T}z}{r^{\mathsf T}r},\qquad
 \sigma_I^2=\frac1{r^{\mathsf T}r}.
\end{equation}
已知分量系数需同时作用于估计和标准不确定度。参考晶畴的全局比例只改变初始强度尺度，不应改变相位搜索；新版的尺度协变更新正是为此设计的。

这些带符号初始观测只负责生成模型。每个模型随后在全部 work 像素组的共同似然上精修，包括弱点、部分重叠和完全重叠。选择分支比较同一批 selection 组。不可辨识的单点分量不需要先分拆，也不会被导出成一张虚构的独立强度表。

输入已经有合理模型时，控制器可以跳过从头相位阶段，执行同一套分组、联合精修、模型选择和审计。没有模型、也没有足够的隔离相位数据时，当前实现会清楚报告缺少可用的初始化证据。这比从任意孪晶总强度构造一个未经物理证明的振幅集合更可检查。

\section{MOF 的理想结构能提供什么}
已知金属节点、配体、化学计量或拓扑，可以明显减少模型空间。例如已知配体的键长角度约束、刚性片段、对称性候选，都可以降低弱数据下的参数方差。一个常见的数学形式是
\begin{equation}
 Q_\lambda(\theta)=\chi^2_{\rm diffraction}(\theta)
                  +\lambda\|g(\theta)-g_0\|_{\Sigma_g^{-1}}^2.
\end{equation}
$g_0$ 来自独立化学知识，$\lambda$ 表示对它的信任程度。数据弱时先验的作用更大，因此结论也更依赖先验。应比较合理的备选模型，例如完整节点与缺陷节点、不同客体占有率、合理的无序形式，而不是只允许朝一个理想结构移动。

当前原生内核已提供距离、角度、相似距离、平面性、ADP 相似和刚性键等限制，以及精确的占有率和位置约束。它没有内置一个经过全面验证的 MOF 拓扑生成器，也不会自动把化学上“看起来顺眼”的原子补成实验已经观察到的原子。理想模型的可信部分用于提出候选，衍射似然负责检验候选。

\section{何时回到上游}
残差的形状可以提示下一个可检验的问题。反射位置随转角系统漂移，提示检查几何；峰两侧残差符号相反，提示质心或宽度；一个区域出现与第二晶格一致的额外峰，提示新的晶畴假设；残差随入射／出射方向变化，提示吸收模型。每次只修改相关前向假设，并保存前后模型。

本版保留的 DIALS 工作流与新的 \code{native.run} 算子可放入同一个持久化有向图。进程具有时间、重试、日志和取消预算；每次尝试保留独立目录、输入快照和结果。重启后可以继续已存在任务，而不是丢掉前面的证据重算。

原生控制器目前自动执行的是明确的候选搜索、局部精修和验证规则。根据所有残差自主重新标定仪器、发现任意混合相、重建未知 MOF 化学的完整闭环，还需要专门的假设生成器与实验标定。接口已经能表达这些分支，但不能把一个工具调用框架写成全部科学决策都已经实现。

\section{判断闭环是否有价值}
对固定审计集，比较新旧流程的预测似然、误差条覆盖、结构偏差和资源消耗。训练残差下降只说明优化器更擅长拟合训练数据。若新增结构先验降低训练误差、却恶化审计似然，控制器应保留这一失败证据。

本次集成测试还做了一个直接检查：将封存审计数据的强度和不确定度极端改变，候选生成、候选选择及最终拟合模型仍应完全不变，只有审计指标改变。这个测试检验的是计算路径是否泄漏，并不替代跨样品的科学验证。

一次最终评价由每次原生求解流程控制。传输失败后以相同输入、请求和随机种子重算，是计算恢复；看到审计结果后更换假设或调参，是新的模型选择。后者需要新的外层评价设计。任务哈希和断点恢复能保证计算来源，不能使已经公开的审计数据重新变得独立。
